\unitopen{20}
\begin{candos}
\item odmienić czasownik zwrotny (\de{sich freuen, sich treffen, sich interessieren});
\item opowiedzieć o hobby i o tym, jak spędzasz wolny czas;
\item użyć \de{sich interessieren für} i \de{sich freuen auf} z poprawnym przyimkiem;
\item umówić się na aktywność (\de{Wollen wir uns treffen?}).
\end{candos}

\section{Słownictwo}
\vocabtable{20}

\section{Zaimek zwrotny \de{sich}}
\begin{center}\small
\begin{tabular}{@{}lll@{}}
\toprule
\textsf{\textbf{Zaimek}} & \textsf{\textbf{Zwrotny (biernik)}} & \textsf{\textbf{Przykład}}\\
\midrule
\de{ich} & \de{mich} & \de{Ich freue mich.}\\
\de{du} & \de{dich} & \de{Du interessierst dich für Musik.}\\
\de{er/sie/es} & \de{sich} & \de{Er trifft sich mit Anna.}\\
\de{wir} & \de{uns} & \de{Wir treffen uns um acht.}\\
\de{ihr} & \de{euch} & \de{Ihr freut euch.}\\
\de{sie/Sie} & \de{sich} & \de{Sie entspannen sich.}\\
\bottomrule
\end{tabular}
\end{center}
Czasowniki zwrotne mają zwykle polskie odpowiedniki z „się”, ale nie zawsze: \de{sich freuen} (cieszyć się), \de{sich
treffen} (spotykać się), ale też \de{sich vorstellen} (przedstawiać się \emph{i} wyobrażać sobie). Połączenia stałe z przyimkami:
\de{sich freuen auf} + biernik (cieszyć się na coś w przyszłości), \de{sich freuen über} + biernik (cieszyć się z czegoś), \de{sich interessieren
für} + biernik, \de{sich erinnern an} + biernik.

\begin{zasada}[Szyk]
Zaimek zwrotny stoi zaraz po odmienionym czasowniku (w zdaniu głównym): \de{Ich treffe \textbf{mich} morgen mit Anna.}
W pytaniu: \de{Treffen wir \textbf{uns} um acht?} W zdaniu podrzędnym przed czasownikiem na końcu: \de{\dots, weil ich \textbf{mich}
auf das Wochenende \textbf{freue}.}
\end{zasada}

\begin{falle}
W polskim „spotkałem Annę” to czasownik przechodni, a po niemiecku \de{Ich treffe Anna} (biernik) \emph{bez} \de{sich}.
Zwrotny dopiero przy wzajemności: \de{Wir treffen \textbf{uns}.} Błąd: \de{Ich interessiere mich in Musik}
(przyimek to \de{für}). Również \de{Ich freue mich für das Wochenende} jest błędem: \de{auf}.
\end{falle}

\begin{szwajcaria}
Szwajcarzy są aktywni w \de{Vereine}: kluby sportowe, śpiewacze i strzeleckie zbierają znajomości, a członkostwo w klubie
to często najszybsza droga do kontaktów. Popularne: \de{Skifahren}, \de{Wandern}, \de{Schwimmen im See}, \de{Velo fahren}.
Zaprosić kogoś na \de{Feierabendbier} to typowy gest zespołu po pracy.
\end{szwajcaria}

\begin{praca}
\de{Ich interessiere mich für Cloud-Architektur.} \de{Ich freue mich auf das neue Projekt.} \de{Wir treffen uns
um zehn im Büro.} \de{Das Team trifft sich jeden Montag.} Na rozmowie: \de{Ich interessiere mich besonders für
\dots} brzmi dużo lepiej niż \de{Ich mag \dots}.
\end{praca}

\section{Dialogi}
\textbf{Dialog 1 --- plany na weekend}
\begin{dialog}
\turn{A}{Was machst du am Wochenende?}{Co robisz w weekend?}
\turn{B}{Ich treffe mich mit Freunden. Wir gehen wandern. Ich freue mich schon darauf.}{Spotykam się z przyjaciółmi. Idziemy w góry. Już się na to cieszę.}
\turn{A}{Wandert ihr oft?}{Czy często wędrujecie?}
\turn{B}{Ja, ungefähr zweimal im Monat. Interessierst du dich auch für die Berge?}{Tak, mniej więcej dwa razy w miesiącu. Czy ty też interesujesz się górami?}
\turn{A}{Ein bisschen. Ich fahre lieber Ski.}{Trochę. Wolę jeździć na nartach.}
\end{dialog}

\section{Ćwiczenia}
\begin{cwiczenia}
\exitem{Uzupełnij: \de{Ich freue \dots{} auf das Wochenende.}}{\de{mich}}
\exitem{Uzupełnij: \de{Wir treffen \dots{} um acht.}}{\de{uns}}
\exitem{Uzupełnij: \de{Du interessierst \dots{} für Musik.}}{\de{dich}}
\exitem{Uzupełnij: \de{Er erinnert \dots{} an den Termin.}}{\de{sich}}
\exitem{Uzupełnij przyimek: \de{Ich interessiere mich \dots{} Sport.}}{\de{für}}
\exitem{Uzupełnij przyimek: \de{Wir freuen uns \dots{} den Urlaub.} (przyszłość)}{\de{auf}}
\exitem{Uzupełnij przyimek: \de{Ich freue mich \dots{} das Geschenk.} (cieszę się z)}{\de{über}}
\exitem{Popraw: \de{Ich treffe mich Anna.}}{\de{Ich treffe Anna.} lub \de{Ich treffe mich mit Anna.}}
\exitem{Popraw: \de{Ich interessiere mich in Musik.}}{\de{Ich interessiere mich für Musik.}}
\exitem{Zaproś kolegę na spotkanie o dwudziestej.}{\de{Treffen wir uns um zwanzig Uhr?}}
\exitem{Przetłumacz: „Cieszę się na nowy projekt.”}{\de{Ich freue mich auf das neue Projekt.}}
\exitem{Przetłumacz: „Czym się interesujesz?”}{\de{Wofür interessierst du dich?}}
\end{cwiczenia}

\weekend{Opowiedz o trzech swoich hobby w osiemnastu zdaniach (po sześć): co robisz, jak często, z kim, czego się spodziewasz.
Użyj \de{sich interessieren für} i \de{sich freuen auf}. Nagraj 90 sekund.}
